\header{Tool Sequence Patterns in Successful vs Unsuccessful Repairs} To understand how repair strategies differ between successful and unsuccessful attempts, we analyze tool sequence patterns using a sliding window approach with $w=3$. We extract consecutive 3-step tool invocations from each trajectory and count pattern frequencies for the two lowest-performing agents, \qwencode and \geminicli. Figures~\ref{fig:qwen_pass_fail_sequences} and \ref{fig:gemini_pass_fail_sequences} show the top five patterns separated by repair outcome (abbreviated as WR=WRITE, RD=READ, TST=TEST, BLD=BUILD, NAV=NAVIGATE, SC=SEARCH\_CONTENT).

% \paragraph{Qwen Code: Structural Navigation Overhead}

\begin{figure}[t]
\centering
\includegraphics[width=0.95\textwidth]{diagrams/tool-sequences/qwen_pass_fail_sequences.png}
\caption{Tool sequence patterns for \qwencode.}
\label{fig:qwen_pass_fail_sequences}
\end{figure}

\qwencode exhibits navigation overhead in both successful and unsuccessful repairs. The most frequent pattern is NAVIGATE $\rightarrow$ TEST $\rightarrow$ NAVIGATE (4.13\% successful, 2.93\% unsuccessful). NAVIGATE appears in 3/5 successful patterns and 4/5 unsuccessful patterns. WRITE $\rightarrow$ NAVIGATE $\rightarrow$ BUILD occurs at 2.64\% in successful repairs and 2.86\% in unsuccessful repairs. Navigation overhead persists regardless of outcome.

% \paragraph{Gemini CLI: Behavioral Shift Toward Over-Modification}

\begin{figure}[t]
\centering
\includegraphics[width=0.95\textwidth]{diagrams/tool-sequences/gemini_pass_fail_sequences.png}
\caption{Tool sequence patterns for \geminicli.}
\label{fig:gemini_pass_fail_sequences}
\end{figure}

\geminicli exhibits different patterns in successful versus unsuccessful repairs. Successful repairs use WRITE $\rightarrow$ BUILD $\rightarrow$ TEST most frequently (9.09\%). Unsuccessful repairs use WRITE $\rightarrow$ WRITE $\rightarrow$ WRITE most frequently (10.37\%), a 2.24-fold increase over successful repairs (4.63\%). WRITE $\rightarrow$ READ $\rightarrow$ WRITE occurs at 6.06\% in unsuccessful repairs. Unsuccessful repairs accumulate consecutive modifications without intermediate compilation or testing.


\qwencode and \geminicli exhibit distinct failure modes. \qwencode shows navigation overhead in both successful and unsuccessful repairs. \geminicli shows a behavioral shift: successful repairs use WRITE $\rightarrow$ BUILD $\rightarrow$ TEST (9.09\%), while unsuccessful repairs use WRITE $\rightarrow$ WRITE $\rightarrow$ WRITE (10.37\%). \qwencode requires architectural changes to reduce navigation overhead. \geminicli requires strategies to prevent consecutive modifications without validation.
